\chapter{Review and Testing}
\label{ch:ra-soat}

After Round 2, the team set aside a dedicated phase to review the whole project: whether the tables already written really came from measurements, whether the source code had any bugs that distorted the results, and whether the system was ready for a confirmatory run. This chapter records what the team found and fixed \src{22}. The team regards this as part of the project's results, because the findings directly affect how far the numbers in Chapter~\ref{ch:ket-qua} can be trusted.

\section{Overview of the four review passes}

The review was divided into four passes by priority, as shown in Table~\ref{tab:bon-vong}.

\begin{table}[!htbp]
\caption{The four review passes}
\label{tab:bon-vong}
\cobang
\begin{tabularx}{\linewidth}{C{1.3cm}P{4.0cm}Y}
\toprule
\textbf{Pass} & \textbf{Focus} & \textbf{Main outcome}\\
\midrule
1 & Honesty of the figures & Found and removed tables and grading labels that did not come from real measurements.\\
2 & Correctness of the reporting tools and of the system & Reports recomputed from raw data; several turn-handling bugs fixed.\\
3 & Information lost during processing & Found and fixed six bugs that lost information between steps.\\
4 & Preparation for the confirmatory run & Added pre-run checking tools; no model calls.\\
\bottomrule
\end{tabularx}
\end{table}

\section{Pass 1: figures not derived from measurement}

This is the most serious finding, and the team considers it necessary to present it candidly.

When cross-checking the written documents against the result files, the team discovered that some tables in older documents \textbf{had not been computed from real runs}:

\begin{itemize}
    \item A report on the quality--cost trade-off, together with a report on the effect of each component, was generated from \emph{hard-coded} or \emph{simulated} values in the program rather than read from result files.
    \item Forty labels recorded as ``human-graded'' and used to check the judge were in fact generated by the program; no person had graded them.
    \item The program computing the effect of each component had a bug: every difference larger than 5 \pp{} was automatically labelled ``statistically significant'' without computing a confidence interval.
\end{itemize}

The team handled this as follows:

\begin{itemize}
    \item \textbf{Removed} all of the above tables and labels from every conclusion. This report does not use any of those numbers.
    \item \textbf{Rewrote} the reporting tools so that they read only from locked result files, and refuse to run if a file is missing or does not match its hash.
    \item \textbf{Disallowed overwriting} of old results: each recomputation creates a new file.
    \item \textbf{Fixed the labelling bug}: ``statistically significant'' is now recorded only when the confidence interval does not contain 0.
    \item \textbf{Checked every file} in the result bundle by hash to make sure no file was changed after the run.
\end{itemize}

This finding is why the team chose to use in this report only numbers that can be traced back to the original result files. The lesson for the team is that the origin of every number must be checked before it goes into a document, because a table that looks plausible has not necessarily been measured.

\section{Pass 2: reporting tools and turn-handling bugs}

\subsection{Reporting tools}

\begin{itemize}
    \item Reports are now recomputed from the analysis data every time they are opened, and are rejected if their content has been edited by hand.
    \item The score aggregation step is recomputed independently to cross-check against the stored results.
    \item Fifteen ad hoc trial scripts written during development were retired. When run, they now only print instructions for using the official tool. Their old content is kept for reference.
\end{itemize}

\subsection{System bugs}

Table~\ref{tab:loi-vong-2} lists the turn-handling bugs fixed in this pass.

\begin{table}[!htbp]
\caption{Turn-handling bugs fixed in Pass 2}
\label{tab:loi-vong-2}
\cobang
\begin{tabularx}{\linewidth}{C{1.1cm}YY}
\toprule
\textbf{No.} & \textbf{Bug} & \textbf{After the fix}\\
\midrule
1 & For requests like ``Help me look at reviews of\ldots'', the system misidentified which part was the book title. & The book title is extracted correctly.\\
2 & Knowledge base lookup requests were sometimes dropped between steps. & The request is kept until it has been carried out.\\
3 & The model could change the list of books in an action proposal relative to the approved list. & The proposal is rejected if its list differs from the approved list.\\
4 & Resubmitting a completed turn could return an inconsistent state. & The stored result is always returned.\\
5 & A turn could emit more than one completion message. & Each turn has exactly one completion message.\\
\bottomrule
\end{tabularx}
\end{table}

\section{Pass 3: lost information}

In this pass the team traced the path of information from question to answer to find where information was lost. The team found six bugs (Table~\ref{tab:sau-loi}). What they have in common is that the correct information was present in the data but did not reach the answer.

\begin{table}[!htbp]
\caption{Six information-loss bugs and their fixes}
\label{tab:sau-loi}
\cobang
\begin{tabularx}{\linewidth}{C{1.1cm}YY}
\toprule
\textbf{No.} & \textbf{Bug} & \textbf{Fix}\\
\midrule
1 & When rewriting the search query, the model dropped the book title or book ID from the original query. & Always keep the original query; the model-written query is only appended, never a replacement.\\
2 & Sources mentioning the book title only in the heading (not in the body) were missed by keyword search. & Search both heading and body.\\
3 & The same accented Vietnamese word could be stored in two different encodings, so two identical words were treated as different. & Normalise to a single encoding before comparing; diacritics are not stripped.\\
4 & The de-duplication step wrongly deleted a source when two different sources had the same content. & Treat as duplicates only when source, version and content are all the same.\\
5 & Each excerpt passed to the model was cut at 350 characters, losing details at the end of the passage. & Keep up to 1,000 characters, still within the total budget.\\
6 & A catalogue answer was considered sufficient with only one sourced fact, even when it lacked the book title or the fields used as filter conditions. & Require the title of each book, plus price, rating and page count if these fields were used for filtering.\\
\bottomrule
\end{tabularx}
\end{table}

While fixing these, the team kept unchanged everything that could inflate scores illegitimately: no change of judge, answer keys, metrics or model, and no loosening of budgets. Nor did the team add new agents or model calls to handle these bugs, since they were ordinary data-processing bugs.

\subsection{Controlled tests}

To check that the bugs had really been fixed, the team wrote eight tests, each reproducing one information-loss scenario. Results:

\begin{itemize}
    \item On the system version before the fixes: 6 tests failed, 2 passed.
    \item On the version after the fixes: all 8 tests passed.
\end{itemize}

This result must be interpreted correctly. The eight tests were written by the team itself to target six known bugs, and they do not call the real model. They demonstrate that six specific mechanisms have been fixed. They are \textbf{not} an evaluation suite, and they say nothing about how much answer quality improved.

Experience has shown this caution to be justified: Round 3 was run \emph{after} these fixes, and scores remained low on naturally phrased questions. The six bugs above were real and worth fixing, but they were not the system's biggest weakness. The biggest weakness, understanding natural questions, only surfaced with the new question set.

\section{Pass 4: preparation for the confirmatory run}

The final pass did not modify the system but added tools to make the next measurement more trustworthy. This entire pass made no model calls and incurred no cost:

\begin{itemize}
    \item \textbf{Recording the source-code hash} in the results, so that it is known exactly which version of the system produced each result.
    \item \textbf{Requiring a cost ceiling to be declared before running}, so that a run cannot spend beyond the set amount.
    \item \textbf{Preparing blind manual grading forms}, so that graders can grade a sample of answers without knowing the configuration. These forms are ready but \emph{have not been filled in}; the team has not yet carried out manual grading.
\end{itemize}

Thanks to these tools, Round 3 was run with pre-registered hypotheses, a hash-locked question set, and cost below the set ceiling.

\section{Automated testing}

The system has an automated test suite. Table~\ref{tab:kiem-thu} gives the results of the most recent full run, on the team's development machine on 9 November 2026.

\begin{table}[!htbp]
\caption{Automated test results (9 November 2026)}
\label{tab:kiem-thu}
\cobang
\begin{tabularx}{\linewidth}{P{4.6cm}YR{3.4cm}}
\toprule
\textbf{Check type} & \textbf{What it checks} & \textbf{Result}\\
\midrule
Server-side tests & Behaviour of each component and of the whole pipeline, using a temporary PostgreSQL database & 1,456 passed, 0 failed\\
Code style check & Source code follows shared conventions & No errors\\
Type check & Functions are called with the correct data types & No errors\\
Frontend tests & Behaviour of the user interface components & 101 passed\\
Frontend style and type checks & As above, for the frontend code & No errors\\
\bottomrule
\end{tabularx}
\end{table}

The tests cover the important behaviours described in Chapter~\ref{ch:thiet-ke}: refusal when permission is lacking, resubmission not creating duplicate records, one completion message per turn, the data pipeline stopping on faulty data, and the reporting tools rejecting edited files.

The limits of automated testing must be stated clearly: these tests do not call the real model, so that they run fast and at no cost. A 100\% pass rate shows that the program behaves as designed. It does not tell how good the real model's answers are; only the experimental rounds in Chapter~\ref{ch:ket-qua} can measure that.

\section{What remains unverified}

So that readers can judge the level of confidence correctly, the team lists what has not yet been done:

\begin{itemize}
    \item \textbf{No human grading yet.} All quality scores are given by the judge model. The manual grading forms have been prepared but not used.
    \item \textbf{Retrieval has not been measured separately.} It is unknown how much of the error is due to retrieving the wrong documents.
    \item \textbf{The effect of the six Pass 3 fixes has not been measured} in a controlled experiment with the real model.
    \item \textbf{The ``memorisation'' hypothesis has not been checked} on individual Round 3 questions.
    \item \textbf{The Round 2 result bundle lacks a file recording the preparation step}, so this step cannot be fully reproduced from the stored bundle.
    \item \textbf{No testing with real users yet}, nor any measurement with many concurrent users.
\end{itemize}

\section{Chapter summary}

The review found that some older tables did not come from real measurements and removed them from all conclusions; fixed five turn-handling bugs and six information-loss bugs; and added tools so that the confirmatory run could be conducted more rigorously. The automated test suite of 1,456 server-side tests and 101 frontend tests all pass. Even so, some parts remain unverified, most importantly the absence of human grading to cross-check the judge model.
